% Part: first-order-logic
% Chapter: introduction
% Section: models-theories

\documentclass[../../../include/open-logic-section]{subfiles}

\begin{document}

\olfileid{fol}{int}{mod}

\olsection{प्रतिमाने आणि उपपत्ती}

प्रथम-क्रम तर्कशास्त्राची विन्यासमीमांसा आणि चिन्हार्थमीमांसा परिभाषित
केल्यानंतर !!{structure}s चे गुणधर्म आणि चिन्हार्थविषयक संकल्पना यांचा
अभ्यास सुरू करता येतो. !!{derivation} पद्धतीही परिभाषित करून त्यांचा
अभ्यास करता येतो. !!{sentence}s चा एखादा संच घेतल्यावर आपण विचारू
शकतो: त्या संचातील सर्व !!{sentence}s कोणती !!{structure}s सत्य करतात?
!!{sentence}s चा संच~$\Gamma$ दिला असता, त्यांची पूर्ति करणाऱ्या
!!a{structure}~$\Struct{M}$ ला \emph{$\Gamma$ चे प्रतिमान} म्हणतात.
~$\Gamma$ पासून सुरुवात करून आपण त्याची प्रतिमाने शोधू शकतो---ती कशी
दिसतात? ती किती मोठी किंवा लहान असावी लागतात? पण एखादे एक
!!{structure} किंवा !!{structure}s चा संग्रह घेऊनही विचारता येते: त्यांत
कोणती !!{sentence}s सत्य आहेत? या !!{structure}s मध्ये ती आणि फक्त तीच
सत्य असतात या अर्थाने त्यांचे \emph{अभिलक्षण} करणारी !!{sentence}s आहेत
का? अशा प्रकारचे प्रश्न \emph{प्रतिमान उपपत्ती} या क्षेत्रातील आहेत.
त्यांवरच \emph{स्वयंसिद्धकीय पद्धत} ही आधारित असते: एखाद्या
!!{structure}s च्या संग्रहाचे उपपत्तीच्या स्वयंसिद्धकांचा, म्हणजे
!!{sentence}s च्या संचाचा, वापर करून वर्णन करणे. स्वयंसिद्धकांपासून
प्रथम-क्रम तर्कशास्त्रात तार्किकरीत्या निष्पन्न होणारी नेमकी तीच
!!{sentence}s स्वयंसिद्धकांच्या सर्व प्रतिमानांत सत्य असतात, या
निरीक्षणामुळे हे शक्य होते.

अतिशय सोपे उदाहरण म्हणून पूर्वक्रमांचा विचार करा. एखाद्या संच~$A$ वरील
परावर्ती आणि संक्रमक असा संबंध~$R$ म्हणजे पूर्वक्रम होय. द्विस्थानी
संबंध $R \subseteq A \times A$ असलेला संच~$A$ म्हणजे एकच द्विस्थानी
संबंधचिन्ह~$\Obj P$ असलेल्या प्रथम-क्रम भाषेला !!a{structure} देण्यासाठी
नेमके जे आवश्यक आहे ते: आपण $\Domain{M} = A$ आणि
$\Assign{\Obj P}{M} = R$ ठेवू. $R$ हा पूर्वक्रम असल्यामुळे तो परावर्ती
आणि संक्रमक आहे आणि हे सांगणारा प्रथम-क्रम तर्कशास्त्रातील
!!{sentence}s चा संच~$\Gamma$ आपल्याला मिळू शकतो:
\begin{align*}
  & \lforall[\Obj v_0][\Atom{\Obj P}{\Obj v_0,\Obj v_0}]\\
  & \lforall[\Obj v_0][\lforall[\Obj v_1][\lforall[\Obj v_2][((\Atom{\Obj P}{\Obj v_0,\Obj v_1} \land \Atom{\Obj P}{\Obj v_1,\Obj v_2}) \lif \Atom{\Obj P}{\Obj v_0,\Obj v_2})]]]
\end{align*}
ही !!{sentence}s म्हणजे फक्त ``कोणत्याही~$x$ साठी $Rxx$'' ($R$ हा
परावर्ती आहे) आणि ``जेव्हा $Rxy$ आणि $Ryz$, तेव्हा $Rxz$ ही'' ($R$ हा
संक्रमक आहे) यांची चिन्हांकने आहेत. !!a{structure}~$\Struct{M}$ हे या
दोन !!{sentence}s~$\Gamma$ चे प्रतिमान तेव्हा आणि केवळ तेव्हाच असते,
जेव्हा $R$ (म्हणजे $\Assign{\Obj P}{M}$) हा~$A$ (म्हणजे
$\Domain{M}$) वरील पूर्वक्रम असतो. दुसऱ्या शब्दांत, $\Gamma$ ची प्रतिमाने
म्हणजे नेमके पूर्वक्रम होत. फक्त~$\Obj P$ हे !!{predicate} असलेल्या
प्रथम-क्रम भाषेत व्यक्त करता येणारा सर्व पूर्वक्रमांचा कोणताही गुणधर्म
(वरील परावर्तिता आणि संक्रमणता यांसारखा) हा~$\Gamma$ मधील दोन
!!{sentence}s पासून तार्किकरीत्या निष्पन्न होतो आणि उलटही तसेच होते.
म्हणून~$\Gamma$ च्या प्रतिमानांविषयी आपण जे काही सिद्ध करू शकतो, ते सर्व
पूर्वक्रमांविषयी सिद्ध केलेले असते.

कोणत्याही विशिष्ट उपपत्तीसाठी आणि प्रतिमानांच्या वर्गासाठी (जसे
$\Gamma$ आणि सर्व पूर्वक्रम) संबंधित प्रथम-क्रम भाषेत काय व्यक्त करता
येते आणि काय व्यक्त करता येत नाही याविषयी रोचक प्रश्न असतील. सर्व भाषा
आणि प्रतिमानवर्गांसाठी रोचक असणारे !!{structure}s चे काही गुणधर्म आहेत;
विशेषतः !!{domain} च्या आकारमानाशी संबंधित गुणधर्म. उदाहरणार्थ,
कोणत्याही~$n \in \PosInt$ साठी !!{domain} मध्ये नेमके
$n$~!!{element}s आहेत असे नेहमी व्यक्त करता येते. अनंतपणे अनेक
!!{sentence}s चा संच वापरून !!{domain} अनंत आहे असेही व्यक्त करता येते.
पण प्रांत सांत आहे किंवा प्रांत !!{nonenumerable} आहे असे व्यक्त करता
येत नाही. प्रथम-क्रम भाषांच्या मर्यादांविषयीचे हे परिणाम संहतता आणि
लोव्हेनहाइम--स्कोलेम प्रमेयांचे परिणाम आहेत.

\end{document}
